% Source: https://artofproblemsolving.com/community/c3h1447823_algebra_aaaaa
% Stable identity: aops-topic-1447823-q01
\textbf{(a)} $\left(x-\dfrac{1}{x}\right)^2=\left(x+\dfrac{1}{x}\right)^2-4=9-4=5$, so $x-\dfrac{1}{x}=\boxed{\pm\sqrt5}$. Both signs occur: $x=\dfrac{3\pm\sqrt5}{2}$, and the two values of $x$ are reciprocals of each other.

\textbf{(b)} $x^2-\dfrac{1}{x^2}=\left(x+\dfrac{1}{x}\right)\left(x-\dfrac{1}{x}\right)=\boxed{\pm3\sqrt5}$ (same sign as in (a)).

\textbf{(c)} $x^2+\dfrac{1}{x^2}=9-2=7$, so
\[
x^3-\dfrac{1}{x^3}=\left(x-\dfrac{1}{x}\right)\left(x^2+1+\dfrac{1}{x^2}\right)=\pm\sqrt5\cdot8=\boxed{\pm8\sqrt5}
\]
(same sign as in (a)). For $x=\dfrac{3+\sqrt5}{2}$ all three values are positive.
